
\SGASubsection{5.8}{End of the proof of \SGARefLink{sga3:xvii:theorem:5:1:1:item:i}{5.1.1(i)}}
\SGARefTarget{sga3:xvii:subsection:5:8}

It remains to prove (i)\SGARefLink{sga3:xvii:appendix:C}{(c)} and (i)(d).

\SGAProof{Proof of \SGARefLink{sga3:xvii:theorem:5:1:1:item:i}{5.1.1(i)}(d): \(U\) smooth and \(H\) connected}
We shall need the following lemma.

\SGAStatement{Lemma 5.8.1}{With the notation of \SGARefLink{sga3:xvii:theorem:5-1-1}{5.1.1}, suppose that
\(U\) is smooth and \(H\) is radicial, and let \(H_1\) be an
algebraic subgroup of \(H\) which has a lift \(H'_1\) in \(E\).
Then \(H\) has a lift \(H'\) in \(E\) which contains \(H'_1\).}
\SGARefTarget{sga3:xvii:lemma:5:8:1}

\SGAProof{Proof}
By induction on the height of \(H/H_1\), we may suppose that
\(H/H_1\) has height \(1\).  Put
\[
  C'=\operatorname{Centr}_E(H'_1).
\]
We claim that the canonical morphism \(C'\longrightarrow H\) is an
epimorphism.  To prove this, we may suppose that \(k\) is
algebraically closed.  The extension \(E\) is then trivial, by
\SGARefLink{sga3:xvii:theorem:5:1:1:item:i}{5.1.1(i)}\SGARefLink{sga3:xvii:appendix:A}{(a)}.  Let \(H''\) be a lift of \(H\) in \(E\), and let
\(H''_1\) be the inverse image of \(H_1\) in \(H''\).  The groups
\(H'_1\) and \(H''_1\) are two lifts of \(H_1\) in \(E\); hence they
are conjugate by an element of \(U(k)\), since \(k\) is
algebraically closed and \(U\) is smooth, by \SGARefLink{sga3:xvii:theorem:5:1:1:item:ii}{5.1.1(ii)}\SGARefLink{sga3:xvii:section:b}{(b)}.
Consequently, to prove the assertion about \(C'\), it is enough to
prove it for
\[
  C''=\operatorname{Centr}_E(H''_1).
\]
But \(C''\) contains \(H''\), so the assertion is clear.

It also follows from \SGARefLink{sga3:xvii:lemma:5:7:1}{5.7.1} that \(C'\cap U\) is smooth.  We may
therefore replace \(E\) by \(C'\), and thus suppose that \(H'_1\) is
central.  Passing then to the quotient by \(H'_1\), we may suppose
that \(H_1=0\) and that \(H\) has height \(1\).

Since \(U\) is smooth, we have the following exact sequence of
\(p\)-Lie algebras (Appendix II, \SGARefLink{sga3:xvii:proposition:3:2}{3.2}):
\[
  (*)\qquad
  0\longrightarrow\operatorname{Lie}U
   \longrightarrow\operatorname{Lie}E
   \longrightarrow\operatorname{Lie}H
   \longrightarrow0.
\]
In view of Appendix II, \SGARefLink{sga3:ii:subsection:2-2}{2.2}, to say that \(E\) is trivial is
equivalent to saying that \((*)\) is a trivial extension of
\(p\)-Lie algebras.  Suppose that \(H\ne0\), and hence that
\(\operatorname{Lie}H\ne0\).  Suppose also that we have found a
nonzero \(p\)-Lie subalgebra
\[
  \mathfrak h_1\subset\mathfrak h=\operatorname{Lie}H
\]
which lifts to a \(p\)-Lie subalgebra
\(\mathfrak h'_1\subset\operatorname{Lie}E\).  By the cited result,
there is a subgroup \(H_1\) of \(H\), with
\(\operatorname{Lie}H_1=\mathfrak h_1\), and a lift \(H'_1\) of
\(H_1\) in \(E\), with
\(\operatorname{Lie}H'_1=\mathfrak h'_1\).  Applying once more the
reduction just described, we are brought back to the same problem
with \(H\) replaced by \(H/H_1\).  Since \(H\) has height \(1\),
\[
  \operatorname{Lie}(H/H_1)
   =\operatorname{Lie}H/\operatorname{Lie}H_1
\]
(loc.\ cit.), and consequently
\[
  \operatorname{rk}\operatorname{Lie}(H/H_1)
  \leq \operatorname{rk}\operatorname{Lie}H-1.
\]
Thus, by induction on the rank of \(\operatorname{Lie}H\), it is
enough, when \(\mathfrak h\ne0\), to find a nonzero Lie subalgebra
of \(\operatorname{Lie}H\) which lifts to \(\operatorname{Lie}E\).
This follows from the next lemma.

\SGAStatement{\SGAInlineRefTarget{sga3:xvii:lemma:5-8-2}Lemma 5.8.2}{Let \(k\) be a field of characteristic
\(p>0\), and let
\[
  \varphi:\mathfrak g\longrightarrow\mathfrak h
\]
be a surjective morphism of finite-rank \(p\)-Lie algebras over
\(k\).  Then:
\begin{enumerate}
\item[\SGAInlineRefTarget{sga3:xvii:lemma:5-8-2:item:i}(i)] If \(\mathfrak h_{\overline{k}}\) is reductive
  (\SGARefLink{sga3:xvii:lemma:4-2-2}{4.2.2}) and nonzero, there exists a reductive Lie subalgebra
  \(\mathfrak h_1\) of \(\mathfrak g\) whose image in
  \(\mathfrak h\) is nonzero.
\item[\SGAInlineRefTarget{sga3:xvii:lemma:5-8-2:item:ii}(ii)] If \(k\) is perfect and \(u\) is a unipotent element of
  \(\mathfrak h\), that is, if \(u^{(p^n)}=0\) for some \(n\), then
  \(u\) lifts to a unipotent element of \(\mathfrak g\).
\end{enumerate}}
\SGARefTarget{sga3:xvii:lemma:5:8:2}

\SGAProof{Proof}
Let \(X\) be a nonzero element of \(\mathfrak h\) in case \SGARefLink{sga3:xvii:lemma:5-8-2:item:i}{(i)}, and
let \(X=u\) in case \SGARefLink{sga3:xvii:lemma:5-8-2:item:ii}{(ii)}; choose a lift \(X'\) of \(X\) in
\(\mathfrak g\).  The \(p\)-Lie subalgebra \(\mathfrak h'\) of
\(\mathfrak g\) generated by \(X'\) is abelian (Exposé VII).

\smallskip
\noindent
\emph{Case \SGARefLink{sga3:xvii:lemma:5-8-2:item:i}{(i)}.}
The description given in \SGARefLink{sga3:xvii:lemma:4-2-2}{4.2.2} shows that the reductive part
(loc.\ cit.) of \(\mathfrak h'_{\overline{k}}\) is already defined
over \(k\); denote it by \(\mathfrak r'\).  We claim that the image
of \(\mathfrak r'\) in \(\mathfrak h\) is nonzero.  We may suppose
that \(k\) is algebraically closed, so that
\[
  \mathfrak h'=\mathfrak r'\oplus\mathfrak u',
\]
where \(\mathfrak u'\) is the unipotent part of
\(\mathfrak h'\).  If the image of \(\mathfrak r'\) in
\(\mathfrak h\) were zero, then the image of \(\mathfrak h'\) in
\(\mathfrak h\) would be unipotent, and hence zero because
\(\mathfrak h\) is reductive (\SGARefLink{sga3:VIIA:2.4:item:ii}{2.4(ii)}); but by construction that
image contains \(X\), a contradiction.

\smallskip
\noindent
\emph{Case \SGARefLink{sga3:xvii:lemma:5-8-2:item:ii}{(ii)}.}
The same argument applies, with the roles of \(\mathfrak r'\) and
\(\mathfrak u'\) interchanged.

\smallskip
\noindent
\emph{End of the proof of \SGARefLink{sga3:xvii:lemma:5:8:1}{5.8.1}.}
Suppose that \(H\ne0\).  By \SGARefLink{sga3:xvii:lemma:5-8-2}{5.8.2} there exists a nonzero reductive
\(p\)-Lie subalgebra
\(\mathfrak h'_1\subset\operatorname{Lie}E\).  Since
\(\operatorname{Lie}U\) is unipotent, by \SGARefLink{sga3:xvii:proposition:4:3:1:item:i}{4.3(i)}, we necessarily
have
\[
  (\operatorname{Lie}U)\cap\mathfrak h'_1=0.
\]
Thus \(\mathfrak h'_1\) is a lift of a \(p\)-Lie subalgebra of
\(\operatorname{Lie}H\).

\smallskip
\noindent
\emph{End of the proof of \SGARefLink{sga3:xvii:theorem:5:1:1:item:i}{5.1.1(i)}(d).}
By \SGARefLink{sga3:xvii:lemma:5:8:1}{5.8.1} there is a family of algebraic subgroups
\(H'_n\) (\(n\in\mathbb N\)) of \(E\) such that
\(H'_{n+1}\) contains \(H'_n\), and such that \(H'_n\) lifts
\({}_{F^n}H\).  The descending sequence of subgroups
\[
  \operatorname{Centr}_E(H'_n)
\]
is stationary; let \(C\) be its stationary value.  The center \(Z\)
of \(C\) contains \(H'_n\) for every \(n\).  Hence the image of
\(Z\) in \(H\) contains \({}_{F^n}(H)\) for every \(n\), and is
therefore an open subgroup of \(H\) (Exposé VII\textsubscript{A},
\S 4).  It is consequently equal to \(H\), since \(H\) is
connected.  To prove that \(E\) is trivial, we may thus replace
\(E\) by \(Z\), and hence suppose that \(E\) is commutative.  We
shall see in \SGARefLink{sga3:xvii:theorem:7:2:1}{7.2.1} that \(E\) then contains a maximal subgroup \(M\)
of multiplicative type, whose formation commutes with extension of
the base field.  Since \(E_{\overline{k}}\) is a trivial extension,
by \SGARefLink{sga3:xvii:theorem:5:1:1:item:i}{5.1.1(i)}\SGARefLink{sga3:xvii:appendix:A}{(a)}, and \(U\) is unipotent, it is clear that \(M\) is
the unique lift of \(H\) in \(E\).

\SGAProof{Proof of \SGARefLink{sga3:xvii:theorem:5:1:1:item:i}{5.1.1(i)}\SGARefLink{sga3:xvii:appendix:C}{(c)}: \(U\) \(k\)-solvable}
Since \((H/H^0)(\overline{k})\) has order prime to \(p\), duality
immediately shows that there exists an integer \(n\) such that
\({}_nH\) is an étale subgroup and the canonical morphism
\[
  {}_nH\longrightarrow H/H^0
\]
is an epimorphism.  Thus \(H/{}_nH\) is connected.  By \SGARefLink{sga3:xvii:lemma:5:2:3:item:i:d}{5.2.3(d)},
there exists a lift \(H'\) of \({}_nH\) in \(E\).  As at the
beginning of the proof of \SGARefLink{sga3:xvii:lemma:5:8:1}{5.8.1}, one sees that
\[
  C=\operatorname{Centr}_E(H')
\]
is a subgroup of \(E\) such that \(C\longrightarrow H\) is an
epimorphism and \(C\cap U\) is smooth.  Replacing \(E\) by \(C\)
and passing to the quotient by \(H'\), we are reduced to the case
where \(U\) is smooth and \(H\) is connected, namely case \SGARefLink{sga3:xvii:lemma:5-8-2:item:i}{(i)}(d).

This completes the proof of \SGARefLink{sga3:xvii:theorem:5-1-1}{5.1.1}.
